% Copyright (c) 2026 Eugene Joseph M. Ragasa.
% SPDX-License-Identifier: Apache-2.0
% Author: Eugene Joseph M. Ragasa
% Formal proof completed: 2026-10-31
\documentclass[10pt]{article}
\usepackage[T1]{fontenc}
\usepackage{amsmath,amssymb,booktabs,tabularx}
\usepackage[margin=0.75in]{geometry}
\usepackage[hidelinks]{hyperref}
\newcolumntype{Y}{>{\raggedright\arraybackslash}X}
\title{Proof Dossier: Exponential Factorization on a Normalized Product Torus}
\author{Eugene Joseph M. Ragasa}
\date{Formal proof completed: 2026-10-31}
\begin{document}
\raggedright
\maketitle

\section{Identity and status}
\begin{tabularx}{\textwidth}{@{}lY@{}}
\toprule
Statement ID & \texttt{BG-FOURIER-THM-124} \\
Status & proved; no admissions \\
Lean declaration & \path{ProjectKoios.BandGauge.Fourier.normalizedProductTorusAverage_exp_separablePhase} \\
Owning module & \path{ProjectKoios/BandGauge/Fourier/NormalizedProductTorus.lean} \\
Author and license & Eugene Joseph M. Ragasa; \texttt{Apache-2.0} \\
\bottomrule
\end{tabularx}

\section{Exact mathematical statement}
Let $D$ be a finite coordinate type, let $\mathbb T_1=\mathbb R/\mathbb Z$ carry its
standard Haar probability measure $\mu$, and give $\mathbb T_1^D$ the product measure
$\mu^D$. For supplied complex phases $\varphi_d:\mathbb T_1\to\mathbb C$,
\[
 \int_{\mathbb T_1^D}
   \exp\!\left(\sum_{d\in D}\varphi_d(q_d)\right)\,d\mu^D(q)
 =\prod_{d\in D}\int_{\mathbb T_1}\exp(\varphi_d(t))\,d\mu(t).
\]
The proposition includes the carrier and measure convention; they are not inferred from external
$\mathbb T^D$ notation.

\section{Source relation and immutable identity}
The statement is source-neutral finite-product integration selected for KSD's flat-parent
comparators. The orthogonal 2-D protocol at lines 414--435 states a uniform $\mathbb T^2$
comparator with a squared one-dimensional Bessel factor; the orthogonal 3-D protocol at lines
359--381 explicitly says uniform $\mathbb T^3$ integration factorizes and states a cubed factor.
The immutable KSD identity is commit
\texttt{fcfc4d5fa76b08c040a02c97a109eb5f5fc06e48}, tree
\texttt{fe8b47b33fcf0b24c5b28f8799ed220e3f6d7d20}; exact paths, hashes, access mode, and the dirty
owner-worktree limitation are recorded in \path{docs/source-audit.md}.

The Formal base is PR~\#44 merge commit
\texttt{24263f08a5b69fc34be3073bc724a3b0ad4b80d6}, tree
\texttt{2cbef2a05174f5ad491392b40dd4265a751ed932}. The owning-module SHA-256 is
\texttt{a9d7a1fb58123cbeeb8e8474b6e8b54d847d59a962945fe09db3f84a42c746c8}.
No topic commit is claimed before checkpoint authorization.

\section{Assumptions and conventions}
The coordinate type is finite. Every coordinate carrier is Mathlib's \texttt{UnitAddCircle}; every
coordinate measure is its standard total-mass-one Haar measure; and the joint measure is
\texttt{Measure.pi}. The phases are arbitrary exact complex functions. No continuity or separate
integrability assumption appears because Mathlib's product-integral identity gives the equality
for the Bochner-integral convention, including its nonintegrable fallback semantics.

\section{Prior art and contribution boundary}
Finite-product Fubini factorization and $\exp(\sum z_d)=\prod\exp(z_d)$ are standard and carry no
local novelty claim. The Lean proof specializes pinned Mathlib theorem
\texttt{MeasureTheory.integral\_fintype\_prod\_eq\_prod}. The local contribution is the narrow,
audited BandGauge interface that fixes the normalized torus carrier and states the exponential
phase consequence without importing a Bessel, runtime, or scientific claim.

\section{Proof structure}
Pointwise, the complex exponential converts the finite phase sum into the product of coordinate
exponentials. Apply finite-product Fubini to that separable field under the product measure. The
resulting factor for coordinate $d$ is exactly the one-circle average of
$\exp(\varphi_d)$, and the finite product of those factors is the right-hand side.

\section{Lean realization and trusted base}
The proof rewrites with \texttt{Complex.exp\_sum} and invokes the pinned finite-product integral
identity from \path{Mathlib/MeasureTheory/Integral/Pi.lean}, lines 105--112. Toolchain: Lean
\texttt{4.35.0} at \texttt{eb6cafb7290e8bd5f2937522d8d804f58f02bafc}; Mathlib
\texttt{56a005855ea17e0b0958f73ec801ae140b0a1679}.
\begin{verbatim}
#print axioms ProjectKoios.BandGauge.Fourier.
  normalizedProductTorusAverage_exp_separablePhase
-- [propext, Classical.choice, Quot.sound]
\end{verbatim}
No repository-owned axiom or admission occurs.

\section{Evidence and interpretation boundaries}
The theorem does not evaluate any one-circle factor as a Bessel function, bind KSD reduced
coordinates or phase scaling, authenticate a finite mesh or quadrature, bound discretization
error, establish finite-to-continuum convergence, or verify a runtime comparator. It supplies no
material, adequacy, novelty, certificate, or publication conclusion.

\section{AI, automation, and review record}
The Lean proof and dossier were AI-assisted and checked by the maintainer. Kernel checking is not
independent mathematical, scientific, source-fidelity, novelty, or publication review.

\section{Unresolved consumer obligations}
A KSD consumer must separately bind its half-open coordinates to \texttt{UnitAddCircle}, prove its
phase is the declared coordinate sum with the frozen amplitude convention, establish the relevant
one-circle integral, and independently verify any finite-mesh approximation and retained artifact.

\section{Review and revision history}
\begin{tabularx}{\textwidth}{@{}lllY@{}}
\toprule Date & Revision & Process & Outcome and limitation \\
\midrule 2026-10-31 & initial dossier & AI-assisted maintainer review & exact normalized product factorization proved; Bessel and runtime bindings remain open \\
\bottomrule
\end{tabularx}
\end{document}
